%%=============================================================================
%% Geautomatiseerde opbouw van de virtuele testomgeving
%% Aanvulling op Sectie sec:virtuele-simulatie-airmax in Hoofdstuk~\ref{ch:proof-of-concept}.
%% Verwijst naar de scripts in bijlage (map poc-scripts/), die hier gedeeltelijk
%% als codefragment worden getoond via minted.
%%=============================================================================

\section{Geautomatiseerde opbouw van de testomgeving}%
\label{sec:automatisering-poc}

Om de virtuele testomgeving uit Sectie~\ref{sec:virtuele-simulatie-airmax}
reproduceerbaar te maken -- een vereiste voor de wetenschappelijke onderbouwing van
deze bachelorproef (Sectie~\ref{sec:meth-verantwoording}) -- werd de volledige opbouw
van de drie VM's herleid tot één enkel script,
\texttt{build\_poc.sh}\footnote{Volledige broncode in Bijlage~\ref{app:poc-scripts} en in de
digitale bijlagen, map \texttt{poc-scripts/}.}. Dit script vereist geen manuele
installatiestappen: het gebruikt de officiële Ubuntu Server \textit{cloud-image} als
basis en configureert elke VM volledig automatisch via \texttt{cloud-init}, een
industriestandaard-mechanisme voor het onbewaakt provisioneren van VM's dat de
klassieke, manuele ISO-installatiewizard overbodig maakt.

\subsection{Overzicht van het buildscript}%
\label{subsec:automatisering-overzicht}

Figuur~\ref{fig:dumbbell-linkemulator} (Sectie~\ref{subsec:sim-architectuur}) toont de
beoogde topologie. \texttt{build\_poc.sh} realiseert deze in vijf stappen:

\begin{enumerate}
  \item het downloaden en converteren van de Ubuntu Server cloud-image naar een
        VirtualBox VDI-basisschijf (eenmalig; nadien hergebruikt voor elke VM);
  \item het opruimen van een eventuele vorige testopstelling, zodat het script
        idempotent herhaalbaar is;
  \item het aanmaken van de drie VM's (\texttt{poc-server}, \texttt{poc-link-emulator},
        \texttt{poc-client}) met de netwerktopologie uit
        Sectie~\ref{subsec:sim-architectuur};
  \item het wachten tot elke VM zijn \texttt{cloud-init}-provisionering heeft
        afgerond, gedetecteerd via de bereikbaarheid van de SSH-poort;
  \item het weergeven van de verbindingsgegevens en de vervolgcommando's.
\end{enumerate}

Elke VM krijgt twee (of, voor de link-emulator, drie) netwerkadapters: een eerste NAT-
adapter, uitsluitend voor beheer (SSH vanaf de hostmachine via een poort-forward, zodat
geen extra host-only netwerkconfiguratie op de hostmachine nodig is), en één of twee
\textit{Internal Network}-adapters die samen de "radiolink" vormen. Onderstaand
fragment uit \texttt{build\_poc.sh} toont hoe de drie VM's met deze topologie worden
aangemaakt:

\begin{listing}[h]
\begin{minted}{bash}
# create_vm <vmname> <role> <hostname> <mgmt_ssh_hostport> <linkA_net> [<linkB_net>]
create_vm "poc-server"        "server"        "poc-server"        \
    "${SSH_PORT_SERVER}"  "${NET_LINK_A}"
create_vm "poc-link-emulator" "link-emulator" "poc-link-emulator" \
    "${SSH_PORT_LINKEMU}" "${NET_LINK_A}" "${NET_LINK_B}"
create_vm "poc-client"        "client"        "poc-client"        \
    "${SSH_PORT_CLIENT}"  "${NET_LINK_B}"
\end{minted}
\caption[Aanmaak van de drie PoC-VM's]{Aanroep van de \texttt{create\_vm}-functie voor
elk van de drie rollen, met hun respectievelijke netwerktopologie.}
\label{lst:create-vm-calls}
\end{listing}

De functie \texttt{create\_vm} zelf (in \texttt{lib/functions.sh}) kloont de
basisschijf, genereert een \texttt{cloud-init}-\textit{seed}-ISO op maat van de rol
(zie Sectie~\ref{subsec:automatisering-cloudinit}), en roept vervolgens
\texttt{VBoxManage} aan om de VM effectief aan te maken en op te starten:

\begin{listing}[h]
\begin{minted}{bash}
VBoxManage createvm --name "${vmname}" --ostype Ubuntu_64 --register
VBoxManage modifyvm "${vmname}" \
    --memory "${VM_MEMORY_MB}" --cpus "${VM_CPUS}" \
    --nic1 nat --natpf1 "ssh,tcp,127.0.0.1,${mgmt_ssh_hostport},,22" \
    --nic2 intnet --intnet2 "${linkA_net}" \
    --audio none --usb off
# (derde NIC enkel voor de link-emulator, zie functions.sh)

VBoxManage storagectl "${vmname}" --name SATA --add sata --controller IntelAHCI
VBoxManage storageattach "${vmname}" --storagectl SATA \
    --port 0 --device 0 --type hdd    --medium "${vm_disk}"
VBoxManage storageattach "${vmname}" --storagectl SATA \
    --port 1 --device 0 --type dvddrive --medium "${seed_iso}"

VBoxManage startvm "${vmname}" --type headless
\end{minted}
\caption[VBoxManage-aanroepen voor VM-aanmaak]{Kern van de \texttt{create\_vm}-functie:
VM-definitie, netwerkadapters, opslag (schijf + cloud-init-ISO) en opstart.}
\label{lst:vboxmanage-create}
\end{listing}

\subsection{Automatische provisionering via cloud-init}%
\label{subsec:automatisering-cloudinit}

Elke rol (\texttt{server}, \texttt{link-emulator}, \texttt{client}) heeft een eigen
\texttt{cloud-init}-configuratiebestand (\texttt{cloud-init/<rol>-user-data.yaml}) dat
bepaalt welke pakketten geïnstalleerd worden, welke gebruikers aangemaakt worden, en
welke configuratiebestanden weggeschreven worden. Dit vervangt elke vorm van manuele
installatie-interactie volledig.

Op de fileserver-VM installeert \texttt{cloud-init} onder meer een SFTP-gebruiker met
een \texttt{chroot}-omgeving (FR-03, BR-01, BR-02) en een \texttt{iperf3}-server voor
de doorvoersnelheidsmetingen uit REQ-01:

\begin{listing}[h]
\begin{minted}{yaml}
users:
  - name: sftpuser
    lock_passwd: false
    plain_text_passwd: "sftp"
    shell: /usr/sbin/nologin

write_files:
  - path: /etc/ssh/sshd_config.d/99-poc-sftp.conf
    content: |
      Match User sftpuser
          ChrootDirectory /srv/sftp
          ForceCommand internal-sftp
          AllowTcpForwarding no
          X11Forwarding no
\end{minted}
\caption[Cloud-init: SFTP-chroot op de server-VM]{Fragment uit
\texttt{server-user-data.yaml}: een SFTP-gebruiker die volledig binnen
\texttt{/srv/sftp} opgesloten blijft.}
\label{lst:cloudinit-server-sftp}
\end{listing}

Op de link-emulator-VM detecteert een klein hulpscript automatisch welke
netwerkinterface bij welk segment hoort (onafhankelijk van de exacte namen die
VirtualBox toekent), brugt de twee "radiolink"-interfaces tot een transparante
Laag-2-verbinding, en herstelt de \texttt{netfilter}-instellingen die anders de
transparantie van die brug zouden verstoren:

\begin{listing}[h]
\begin{minted}{bash}
# /usr/local/bin/configure-network.sh (rol "bridge")
cat > /etc/netplan/60-poc.yaml <<NETPLAN
network:
  version: 2
  ethernets:
    ${MGMT_IF}:   { dhcp4: true }
    ${LINK_A_IF}: { dhcp4: false }
    ${LINK_B_IF}: { dhcp4: false }
  bridges:
    br0:
      interfaces: [${LINK_A_IF}, ${LINK_B_IF}]
      dhcp4: false
NETPLAN
netplan apply
\end{minted}
\caption[Cloud-init: L2-brug op de link-emulator]{Kern van
\texttt{configure-network.sh} in brug-modus: de twee "radiolink"-interfaces worden
samengevoegd tot één transparant Laag-2-segment, zodat server en client zich
gedragen alsof ze rechtstreeks door de lucht verbonden zijn.}
\label{lst:cloudinit-bridge}
\end{listing}

\subsection{Van scenario naar netem: één commando per meting}%
\label{subsec:automatisering-scenario}

Bovenop de basisopstelling installeert \texttt{cloud-init} op de link-emulator-VM ook
het scenario-script uit Sectie~\ref{subsec:sim-parameters}
(\texttt{poc-set-scenario.sh}), dat bij het opstarten automatisch het "ideaal"-scenario
activeert en nadien via SSH op afstand kan worden gewijzigd met het hostscript
\texttt{scenario-select.sh}:

\begin{listing}[h]
\begin{minted}{bash}
./scenario-select.sh ideaal      # 0--15 km, sterk signaal
./scenario-select.sh marginaal   # 15--30 km, adaptieve modulatie
./scenario-select.sh dfs         # radardetectie: 10s onderbreking + hervatting
\end{minted}
\caption[Scenario wisselen vanaf de host]{Het hostscript \texttt{scenario-select.sh}
schakelt op afstand tussen de drie kanaalscenario's uit
Figuur~\ref{fig:airmax-scenarios}, zonder dat de server- of client-VM herstart moeten
worden.}
\label{lst:scenario-select}
\end{listing}

Het bijhorende hostscript \texttt{run-tests.sh} voert vervolgens, voor elk scenario, de
standaard teststeekproef uit: een \texttt{iperf3}-doorvoertest (REQ-01), een
latency-/jittermeting met \texttt{mtr} (REQ-02), en een effectieve SFTP-overdracht van
een testbestand (FR-01, FR-02), zodat de resultaten in
Hoofdstuk~\ref{ch:resultaten-en-evaluatie} voor elk scenario op identieke wijze en
zonder manuele tussenkomst verzameld worden.

\subsection{Volledige werkwijze}%
\label{subsec:automatisering-workflow}

Samengevat volstaan, na installatie van VirtualBox op de hostmachine, de volgende
commando's om de volledige testopstelling op te bouwen, een scenario te kiezen, en de
metingen uit te voeren:

\begin{listing}[h]
\begin{minted}{bash}
cd poc-scripts/
./build_poc.sh                  # bouwt en start de drie VM's (eenmalig ~10-15 min)
./scenario-select.sh marginaal  # kies het gewenste kanaalscenario
./run-tests.sh                  # voer de meetsteekproef uit en verzamel de resultaten
\end{minted}
\caption[Volledige workflow]{De volledige workflow van installatie tot meting, herleid
tot drie commando's.}
\label{lst:workflow-volledig}
\end{listing}

Deze volledige automatisering draagt bij aan de reproduceerbaarheid van dit onderzoek:
elke meting in Hoofdstuk~\ref{ch:resultaten-en-evaluatie} kan door een derde partij,
met enkel VirtualBox en de meegeleverde scripts, exact worden gereproduceerd zonder
kennis van de onderliggende manuele configuratiestappen te vereisen.
